%% notes-tex/030-scaling-study/sections/2-grid.tex

\section{The grid\secstatus{todo}}
\label{sec:grid}

Every 030 run is the composite closure arm on the query-pair-disjoint split,
\texttt{cgt\_030\_s3\_q\_tok\_embfit\_004}: the per-entry trainer with the dataset
token, the strict pool S3Q, every stored entry a training row, scored on the
pinned 37,705 validation triples. What varies is the size and the training
subset (Figure~\ref{fig:grid}). Every 025 run is the same arm's merged-label twin,
\texttt{cgt\_s3\_q\_kl\_embfit\_037}, on the same subsets at the current size.
Batch is 256 records on one A40, the learning rate 2.5e-4 constant, the
normalization constants those of the full pool on every subset
(\texttt{transforms.fit\_stats\_strict=false}), so the optimization target does
not move with $D$ on 030. The 025 trainer has no such switch and refits its
constants on each subset's training records; the losses are compared in label
units, where that difference does not enter, but the training objectives differ
on that side of the pair.

\tcfig{figures/030-scaling-grid.pdf}{%
The grid as queued. Columns are the nested subsets of S3Q, drawn by query pair
with seed 42 so the 1/64 pairs lie inside the 1/16 pairs inside the 1/4 pairs;
rows are the five sizes with their non-embedding $N$. A filled mark is a queued
run; three marks are seeds 1 to 3, placed where runs are cheapest (the two
smallest subsets) at the two end sizes; the red ring is the learning-rate check
at the smallest size on 1/4. The purple row is the 025 twin at 180x8. The strip
below gives each column's epoch budget, the records seen under it, and the wall
time at the assumed pace; the dashed line is the shared budget of the matched
read. Planned figure drawn from the grid's own tables; no run has reported.
}{fig:grid}
\src{notes/assets/drawio/030-scaling-study.gen.py}

\notesec{The size ladder}
Width and depth scale together in the current configuration's ratio at nine
attention heads (Table~\ref{tab:sizes}): 45x2, 63x3, 90x4, 126x6 and 180x8, the
last being the size every 030 and 025 arm has trained at. $N$ is read from the
model's own tally at each size, the count the compute callback logs, and the
180x8 count equals the one W\&B holds for the instrumented IGB run. The ladder
spans 1.66 decades, under the two decades the explainer's recipe asks for, and
the composite preprocessor that projects the sequence embeddings to the hidden
width is excluded from $N$ as the embedding is.

%% GENERATED FILE -- do not hand-edit.
%% SOURCE: experiments/030-solid-growth-multi/scripts/scaling_grid_table_030.py, reading the model's num_parameters at each size (checked against W&B 53qqijnu)

\begin{table}[htbp]
\centering\footnotesize
\caption{The $N$ ladder. Width and depth scale together at 9 attention heads; $N$ is the non-embedding count the compute callback logs, read from the model's own tally at each size, with the gene table and the composite preprocessor excluded. The 180x8 count equals the one W\&B holds for the instrumented IGB run (53qqijnu, 3,587,350), asserted by the generating script.}
\label{tab:sizes}
\begin{tabular}{rrrrr}
\toprule
width $\times$ depth & head dim & $N$ (non-embedding) & $N / N_{180\times 8}$ & $\log_{10} N$ \\
\midrule
45 $\times$ 2 & 5 & 79,375 & 0.022 & 4.90 \\
63 $\times$ 3 & 7 & 202,612 & 0.056 & 5.31 \\
90 $\times$ 4 & 10 & 509,260 & 0.142 & 5.71 \\
126 $\times$ 6 & 14 & 1,378,444 & 0.384 & 6.14 \\
180 $\times$ 8 & 20 & 3,587,350 & 1.000 & 6.55 \\
\bottomrule
\end{tabular}
\end{table}


\notesec{The data ladder}
Subsets are drawn by the unit that leaks, the query pair, so a smaller subset
removes whole query screens with every triple they reach, and the closure doubles
follow their triples (Table~\ref{tab:subsets}). Singles are in every subset, so
the single-mutant fitness supervision does not vary with $D$, and the pinned
validation and test triples are in every pool unchanged. The same shuffle on the
same sorted pair list gives the same kept pairs on the 025 build, pair for pair,
so a 030 subset and its 025 twin hold the same genotypes and differ only in what
is kept per genotype. $D$ is counted in both units. The ratio between them
drifts from 1.56 rows per record at 1/64 to 2.24 at the full pool, because the
Costanzo and Kuzmin double-mutant screens that supply most of the extra rows
cover the recurring pairs more densely; the 030 curve is therefore fitted in
both units and the two exponents reported side by side.

%% GENERATED FILE -- do not hand-edit.
%% SOURCE: experiments/030-solid-growth-multi/scripts/scaling_grid_table_030.py, reading subset_s3q_fractions_030_summary.json, query_pair_disjoint_splits_030.json.gz and the epoch table of delta_submit_scaling_030.sh

\begin{table}[htbp]
\centering\footnotesize
\caption{The $D$ ladder and its budgets. Subsets are nested by query pair (seed 42, the same kept pairs on the 025 build); singles are in every subset and the pinned validation and test triples are untouched. Records are the 025 unit of $D$; gene\_interaction entry rows are the 030 unit, what the per-entry loss averages over, and rows per record is the ratio between them. Epochs are the submit script's budget, records seen is records $\times$ epochs, and hours is that count at the ASSUMED pace of 68 records/s on one A40 at 180x8 plus stage-in, pending the calibration jobs.}
\label{tab:subsets}
\begin{tabular}{lrrrrrrrrr}
\toprule
subset & pairs & triples & doubles & records & gi rows & rows/rec & epochs & records seen & h at pace \\
\midrule
1/64 & 6 & 5,211 & 10,428 & 21,333 & 33,348 & 1.56 & 240 & 5.12M & 20.9 \\
1/16 & 21 & 21,045 & 42,088 & 68,827 & 131,613 & 1.91 & 120 & 8.26M & 33.8 \\
1/4 & 83 & 79,677 & 159,031 & 244,402 & 509,458 & 2.08 & 40 & 9.78M & 40.0 \\
full & 331 & 301,236 & 739,148 & 1,046,078 & 2,340,771 & 2.24 & 10 & 10.46M & 42.8 \\
\bottomrule
\end{tabular}
\end{table}


\notesec{Budgets}
Every point is one 48 h job, because a dependent job accrues no queue priority
on Delta and a chain pays a fresh queue wait per link. The epoch budget is keyed
by subset: 240, 120, 40 and 10 epochs for 1/64 through the full pool, which is
what 48 h holds at the 180x8 pace for the two largest subsets. That pace, 68
records per second on one A40, is read from the 025 four-GPU runs and is an
assumption until the calibration jobs report; the first version of the submit
script assumed four times that and would have timed out every full-data job
before its tenth epoch. The budget table is keyed by subset only, so a small size
on the full pool gets the same ten epochs as 180x8 whether or not it could have
trained longer, which Section~\ref{sec:support} returns to. The 180x8 full-data
point at ten epochs is the headline point of both the $D$ and $N$ axes and is
not converged: the same arm on IGB was at validation Pearson 0.384 over epochs 13
to 22 and still rising, and the random-split seed took about 150 epochs to turn.

\notesec{What is recorded per run}
The compute callback (\texttt{torchcell.trainers.compute\_accounting}) writes
$N$ at fit start, and at every epoch end the records per epoch, the cumulative
records seen, the measured FLOPs per step and their cumulative sum, and the GPU
hours, together with the epoch index. FLOPs are counted by
\texttt{torch.utils.flop\_counter.FlopCounterMode} around five real training
steps (batch indices 1 to 5, forward and backward), and the mean, the minimum and
the maximum are logged, so a per-entry batch whose row count varies shows its
spread rather than one draw; $6ND$ is a rule for token transformers and does not
describe a gene-token model with graph penalties. The count covers the matmul
family only and is a lower bound, comparable across sizes of one architecture.
The callback's state is saved with the checkpoint, so a chain of jobs continues
its totals.

%% GENERATED FILE -- do not hand-edit.
%% SOURCE: experiments/030-solid-growth-multi/scripts/scaling_grid_table_030.py, reading the loops of delta_submit_scaling_030.sh and delta_submit_scaling.sh

\begin{table}[htbp]
\centering\footnotesize
\caption{The job list, 32 jobs of the 030 grid and 8 of the 025 pair, one GPU each on the Delta A40 partition, 48 h apiece except the two 3 h calibration jobs (400 training batches at 180x8 and 45x2, which report the pace and the measured FLOPs per step). Every 030 job is the composite closure arm on the query-pair-disjoint split (cgt\_030\_s3\_q\_tok\_embfit\_004) at batch 256; every 025 job is its merged-label twin (cgt\_s3\_q\_kl\_embfit\_037) at 180x8.}
\label{tab:jobs}
\begin{tabular}{llrrrr}
\toprule
sweep & what varies & jobs & sizes & subsets & seeds \\
\midrule
030\_sclw1 & calibration & 2 & 45x2, 180x8 & full & 1 \\
030\_sclw1 & size $\times$ subset & 20 & 45x2, 63x3, 90x4, 126x6, 180x8 & 1/64, 1/16, 1/4, full & 1 \\
030\_sclw1 & seed replicates & 8 & 45x2, 180x8 & 1/64, 1/16 & 2, 3 \\
030\_sclw1 & learning-rate check (1e-4, 1e-3) & 2 & 45x2 & 1/4 & 1 \\
025\_sclw1 & subset & 4 & 180x8 & 1/64, 1/16, 1/4, full & 1 \\
025\_sclw1 & seed replicates & 4 & 180x8 & 1/64, 1/16 & 2, 3 \\
\bottomrule
\end{tabular}
\end{table}

